% આ ભાગનો પૂર્ણ સંપાદનીય પાઠ છે. શૈલી, ફોન્ટ, ચિત્રો અને પ્રકરણસંદર્ભ માટે સંપૂર્ણ સ્રોત ZIP તથા reader/full-reader.tex જુઓ. આ એકલી સ્વતંત્ર કમ્પાઇલ થતી મુખ્ય ફાઇલ નથી.

% BEGIN OLP-0055
\olpart{pl}{વિધાનાત્મક તર્કશાસ્ત્ર}
\phantomsection\label{guunit:OLP-0055}


\begin{editorial}
  આ ભાગમાં શાસ્ત્રીય વિધાનાત્મક તર્કશાસ્ત્રની સામગ્રી છે. પ્રથમ પ્રકરણ
  પ્રમાણમાં પ્રાથમિક છે અને તેમાં માત્ર વ્યાખ્યાઓ અને પરિણામોની યાદી છે;
  ઘણી સાબિતીઓ કરી નથી, પણ અભ્યાસપ્રશ્નો તરીકે છોડી છે. પ્રમાણપદ્ધતિઓ
  અને પૂર્ણતા પ્રમેયની સામગ્રી પ્રથમ-ક્રમ તર્કશાસ્ત્રના ભાગમાંથી,
  ``FOL'' ટૅગને false રાખીને, સામેલ કરી છે. તેથી વિધેયો, પદો અને
  પરિમાણકોને લગતું બધું રહી જાય છે અને સંરચનાઓ~$\Struct{M}$ની
  વાતને બદલે સત્યમૂલ્ય-નિયુક્તિઓ~$\pAssign{v}$ની વાત આવે છે.

  આ ભાગને આગળ વધારીને વધુ વિગત ઉમેરવાની અને સત્યતાફલનલક્ષી પૂર્ણતા
  જેવા બીજા વિષયો તથા પરિણામો સામેલ કરવાની યોજના છે.
\end{editorial}

\begingroup
\makeatletter\def\ol@sectioncs{section}\makeatother

% BEGIN OLP-0056
\olchapter{pl}{syn}{વાક્યરચના અને અર્થવિચાર}
\olchapterid{pl}{syn}
\phantomsection\label{guunit:OLP-0056}


\begin{editorial}
  આ માત્ર વ્યાખ્યાઓનો બહુ ટૂંકો સાર છે. સૂત્રો પર અનુમાન વડે થતી
  સાબિતીઓનો સહેલો પરિચય અને ઘણાં વધુ ઉદાહરણો આપીને તેને વિસ્તૃત
  કરવો જોઈએ.
\end{editorial}

\begingroup
\makeatletter\def\ol@sectioncs{section}\makeatother

% BEGIN OLP-0057
\olfileid{pl}{syn}{int}

\olsection{પરિચય}
\phantomsection\label{guunit:OLP-0057}


વિધાનાત્મક તર્કશાસ્ત્રમાં $\lnot$, $\land$, $\lor$, $\lif$ અને
$\liff$ એ વિધાનાત્મક સંયોજકો વડે પ્રતિજ્ઞાપન ચલોમાંથી
બનાવેલાં સૂત્રોનો અભ્યાસ થાય છે. સહજ રીતે કહીએ તો,
પ્રતિજ્ઞાપન ચલ~$p$ સત્ય કે અસત્ય હોઈ શકે એવા વાક્ય
અથવા વિધાનને રજૂ કરે છે. સૂત્રમાંના પ્રતિજ્ઞાપન ચલનું
``સત્યમૂલ્ય'' નક્કી થાય એટલે વિધાનાત્મક સંયોજકો વડે તેમાંથી બનાવેલાં
સૂત્રોનું સત્યમૂલ્ય પણ નક્કી થાય છે. આપણે
વિધાનાત્મક તર્કશાસ્ત્રને \emph{સત્યતાફલનલક્ષી} કહીએ છીએ, કારણ કે તેનો
અર્થવિચાર સત્યમૂલ્યોનાં વિધેયો વડે અપાય છે. ખાસ કરીને, વિધાનાત્મક
તર્કશાસ્ત્રમાં સત્ય અને અસત્યના બીજા કોઈ નિર્ધારણને ધ્યાનમાં લેતા નથી;
દાખલા તરીકે, કોઈ વાત માત્ર સંજોગવશાત્ સત્ય હોવાને બદલે અનિવાર્યપણે
સત્ય છે કે નહિ, તે સત્ય છે એવું જાણીતું છે કે નહિ, અથવા તે હમણાં
સત્ય છે કે પહેલાં સત્ય હતી કે પછી સત્ય થશે. આપણે માત્ર સત્ય
($\True$) અને અસત્ય ($\False$) એ બે સત્યમૂલ્યો લઈએ છીએ; એટલે કોઈ
કથન ન તો સત્ય ન તો અસત્ય, અથવા માત્ર અડધું સત્ય હોઈ શકે એવી શક્યતા
ચર્ચામાંથી બહાર રાખીએ છીએ. વળી, આપણે માત્ર એવા સંયોજકો પર ધ્યાન
આપીએ છીએ જેમનાથી બનેલા સૂત્રનું સત્યમૂલ્ય તેના ભાગોનાં
સત્યમૂલ્યો વડે સંપૂર્ણપણે નક્કી થાય છે (અને, માનો કે, તેના અર્થ વડે
નહિ). ખાસ કરીને, અંગ્રેજી ભાષાનાં શરતી વાક્યોનું સત્યમૂલ્ય આ અર્થમાં
સત્યતાફલનલક્ષી છે કે નહિ તે વિવાદાસ્પદ છે. સત્યમૂલ્યલક્ષી પ્રેરણ
$\lif$ એવું છે; બીજા તર્કશાસ્ત્રો સત્યતાફલનલક્ષી ન હોય એવાં શરતી
વિધાનોનો અભ્યાસ કરે છે.

સત્યતાફલનલક્ષી વિધાનાત્મક તર્કશાસ્ત્રનો સિદ્ધાંત અને અધિસિદ્ધાંત
વિકસાવવા માટે પહેલાં તેના પદોની વાક્યરચના અને અર્થવિચાર વ્યાખ્યાયિત
કરવાં પડે. સંયોજકો વડે પ્રતિજ્ઞાપન ચલોમાંથી
સૂત્રો રચવાની એક રીત આપણે વર્ણવીશું. બીજી વ્યાખ્યાઓ પણ શક્ય
છે. બીજી પદ્ધતિઓ જુદા સંકેતો પસંદ કરશે, સંયોજકોના જુદા ગણને મૂળભૂત
ગણશે અને કૌંસ જુદી રીતે વાપરશે (અથવા કહેવાતી પોલિશ સંજ્ઞાપદ્ધતિની
જેમ જરા પણ નહિ વાપરે). છતાં, બધી રીતોમાં એક વાત સામાન્ય છે: રચનાના
નિયમો સૂત્રોના ગણને \emph{અનુમાનાત્મક રીતે} વ્યાખ્યાયિત કરે છે.
નિયમો યોગ્ય રીતે આપ્યા હોય, તો દરેક પદ રચનાના નિયમો અનુસાર મૂળભૂત રીતે
એક જ રીતે બની શકે. પદોને \emph{એકમાત્ર વાંચી શકાય} એવી અનુમાનાત્મક
વ્યાખ્યા મળવાથી એ જ રીત---અનુમાનાત્મક વ્યાખ્યા---વાપરીને આપણે આ
પદોને અર્થ આપી શકીએ છીએ.

પદોને અર્થ આપવાનું ક્ષેત્ર અર્થવિચાર છે. વિધાનાત્મક તર્કશાસ્ત્રના
અર્થવિચારમાં સત્યમૂલ્ય-નિયુક્તિમાં થતો સંતોષ મુખ્ય ખ્યાલ છે.
સત્યમૂલ્ય-નિયુક્તિ~$\pAssign{v}$ પ્રતિજ્ઞાપન ચલોને સત્યમૂલ્યો
$\True$, $\False$ નિયુક્ત કરે છે. દરેક સત્યમૂલ્ય-નિયુક્તિ, દરેક
સૂત્ર~$!A$ માટે સત્યમૂલ્ય $\pValue{v}(!A)$ નક્કી કરે છે.
$\pValue{v}(!A)=\True$ હોય તો અને તો જ સૂત્ર,
સત્યમૂલ્ય-નિયુક્તિ~$\pAssign{v}$માં સંતોષાય છે---આને આપણે
$\pSat{v}{!A}$ તરીકે લખીએ છીએ. આ સંબંધને પણ $!A$ની રચના પર અનુમાન
વડે વ્યાખ્યાયિત કરી શકાય છે: તાર્કિક સંયોજકોનાં સત્યવિધેયો વાપરીને,
માનો કે, $!A\land !B$નો સંતોષ $!A$ અને $!B$ના સંતોષ (અથવા
અસંતોષ)ના આધારે વ્યાખ્યાયિત થાય છે.

વાક્યો માટેના સંતોષસંબંધ $\pSat{v}{!A}$ને આધારે પછી આપણે પુનરુક્તિ,
ફલિતતા અને સંતોષ્યતાના મૂળ અર્થવિચારી ખ્યાલો વ્યાખ્યાયિત કરી શકીએ
છીએ. દરેક સત્યમૂલ્ય-નિયુક્તિ તેને સંતોષે, એટલે કે દરેક $\pAssign{v}$ માટે
$\pValue{v}(!A)=\True$ હોય, તો સૂત્ર પુનરુક્તિ છે અને આપણે
$\Entails !A$ લખીએ છીએ. જો $\Gamma$નાં બધાં સૂત્રોને સંતોષતું
દરેક સત્યમૂલ્ય-નિયુક્તિ, $!A$ને પણ સંતોષે, તો સૂત્રોનો ગણ તેને
ફલિત કરે છે અને આપણે $\Gamma\Entails !A$ લખીએ છીએ. અને કોઈ એક
સત્યમૂલ્ય-નિયુક્તિ, સૂત્રોના ગણમાંનાં બધાં સૂત્રોને એકસાથે સંતોષે તો
તે ગણ સંતોષ્ય છે. સૂત્રો અનુમાનાત્મક રીતે વ્યાખ્યાયિત છે અને
સંતોષ પણ સૂત્રોની રચના પર અનુમાન વડે વ્યાખ્યાયિત થાય છે; તેથી
અર્થવિચારના ગુણધર્મો સાબિત કરવા અને વ્યાખ્યાયિત અર્થવિચારી ખ્યાલોને
પરસ્પર જોડવા આપણે અનુમાન વાપરી શકીએ છીએ.
% END OLP-0057
\endgroup


\begingroup
\makeatletter\def\ol@sectioncs{section}\makeatother

% BEGIN OLP-0058
\olfileid{pl}{syn}{fml}

\olsection{વિધાનાત્મક સૂત્રો}
\phantomsection\label{guunit:OLP-0058}


વિધાનાત્મક તર્કશાસ્ત્રનાં સૂત્રો,
\emph{પ્રતિજ્ઞાપન ચલો}માંથી,
  વિધાનાત્મક અચળ~$\lfalse$
    અને વિધાનાત્મક અચળ~$\ltrue$ \emph{તાર્કિક સંયોજકો} વડે
રચાય છે.

\begin{enumerate}
\item પ્રતિજ્ઞાપન ચલો $\Obj p_0$, $\Obj p_1$, \dots નો
  ગણનીય ગણ~$\PVar$.
\item અસત્ય માટેનો વિધાનાત્મક અચળ~$\lfalse$.
\item સત્ય માટેનો વિધાનાત્મક અચળ~$\ltrue$.
\item તાર્કિક સંયોજકો:
  \startycommalist
  \ycomma $\lnot$ (નિષેધ)%
  \ycomma $\land$ (સંયોજન)%
  \ycomma $\lor$ (વિયોજન)%
  \ycomma $\lif$ (પ્રેરણ)%
  \ycomma $\liff$ (દ્વિમુખી પ્રેરણ)%
\item વિરામચિહ્નો: (, ) અને અલ્પવિરામ.
\end{enumerate}

વિધાનાત્મક તર્કશાસ્ત્રની આ ભાષાને આપણે $\Lang L_0$થી દર્શાવીએ છીએ.

.




\begin{intro}
ઉપર આપણે વાપર્યાં તેનાથી જુદી પરિભાષા અને સંકેતોથી તમે પરિચિત હોઈ
શકો. તર્કશાસ્ત્રનાં પુસ્તકો (અને શિક્ષકો) ``નિષેધ'' માટે સામાન્ય રીતે
$\sim$, $\neg$ અને~!\; તથા ``સંયોજન'' માટે $\wedge$, $\cdot$ અને
$\&$માંથી કોઈ સંકેત વાપરે છે. ``શરતી વિધાન'' અથવા ``પ્રેરણ'' માટે
સામાન્ય સંકેતો $\rightarrow$, $\Rightarrow$ અને $\supset$ છે.
``દ્વિશરતી વિધાન'', ``દ્વિમુખી પ્રેરણ'' અથવા
  ``(સત્યમૂલ્યલક્ષી) સમતુલ્યતા'' માટે $\leftrightarrow$,
  $\Leftrightarrow$ અને $\equiv$ વપરાય છે.
$\lfalse$ સંકેતને જુદી જગ્યાએ ``અસત્યતા'',
  ``ફાલ્સમ'', ``અસંગતિ'' અથવા ``તળિયું'' કહે છે.
$\ltrue$ સંકેતને જુદી જગ્યાએ ``સત્ય'',
  ``વેરમ'' અથવા ``ટોચ'' કહે છે.
\end{intro}

\begin{defn}[સૂત્ર]
\ollabel{defn:formulas}
વિધાનાત્મક તર્કશાસ્ત્રનાં \emph{સૂત્રો}નો ગણ~$\Frm[L_0]$
નીચે પ્રમાણે અનુમાનાત્મક રીતે વ્યાખ્યાયિત થાય છે:
\begin{enumerate}
\item $\lfalse$ એ આણ્વિક સૂત્ર છે.

\item $\ltrue$ એ આણ્વિક સૂત્ર છે.

\item દરેક પ્રતિજ્ઞાપન ચલ~$\Obj p_i$ આણ્વિક
  સૂત્ર છે.

\item જો $!A$ સૂત્ર હોય, તો $\lnot !A$ પણ
  સૂત્ર છે.

\item જો $!A$ અને $!B$ સૂત્રો હોય, તો $(!A \land
  !B)$ પણ સૂત્ર છે.

\item જો $!A$ અને $!B$ સૂત્રો હોય, તો $(!A \lor !B)$
  પણ સૂત્ર છે.

\item જો $!A$ અને $!B$ સૂત્રો હોય, તો $(!A \lif !B)$
  પણ સૂત્ર છે.

\item જો $!A$ અને $!B$ સૂત્રો હોય, તો $(!A \liff !B)$
  પણ સૂત્ર છે.

\item બીજું કશું સૂત્ર નથી.
\end{enumerate}
\end{defn}

\begin{explain}
સૂત્રોની વ્યાખ્યા \emph{અનુમાનાત્મક વ્યાખ્યા} છે. મૂળભૂત રીતે,
આપણે સૂત્રોનો ગણ અનંત સંખ્યામાં તબક્કાઓમાં રચીએ છીએ. પહેલા
તબક્કામાં બધાં આણ્વિક સૂત્રોને સૂત્ર જાહેર કરીએ છીએ; આ વ્યાખ્યાના
પહેલા થોડા કિસ્સાઓ, એટલે કે
$\ltrue$, %
$\lfalse$, %
$\Obj p_i$ના કિસ્સાઓને અનુરૂપ છે. તેથી ``આણ્વિક સૂત્ર''નો અર્થ
આ સ્વરૂપનું કોઈ પણ સૂત્ર થાય છે.

વ્યાખ્યાના બીજા કિસ્સાઓ અગાઉથી રચેલાં સૂત્રોમાંથી નવાં
સૂત્રો રચવાના નિયમો આપે છે. બીજા તબક્કામાં આણ્વિક
સૂત્રોમાંથી સૂત્રો રચવા આપણે એ નિયમો વાપરી શકીએ છીએ.
ત્રીજા તબક્કામાં આણ્વિક સૂત્રો અને બીજા તબક્કામાં મળેલાં સૂત્રોમાંથી
નવાં સૂત્રો રચીએ છીએ, અને આમ આગળ ચાલે છે. આવા કોઈ તબક્કે છેવટે
રચાતી દરેક વસ્તુ, અને માત્ર એવી વસ્તુ, સૂત્ર છે.
\end{explain}

દ્વિસ્થાની સંયોજક~$\ast$ વડે $!B$, $!C$માંથી રચેલું સૂત્ર
$(!B \ast !C)$ લખતી વખતે આપણે ઘણી વાર સૌથી બહારની કૌંસજોડી છોડી
દઈને માત્ર~$!B \ast !C$ લખીશું.


\sourcecorrection{OLPL-001}{મૂળની \texttt{formulas.tex}માં વ્યાખ્યાયિત
પ્રેરણના વિયોજનવાળા વિકલ્પને $\lnot !A\lor !B)$ લખતાં એક વધારાનું
અંતકૌંસ આવે છે. અહીં અર્થાત્ બનતું સુઘડ સૂત્ર $\lnot !A\lor !B$
લખ્યું છે.}

\begin{defn}[વાક્યરચનાત્મક અભિન્નતા]
$\ident$ સંકેત, સંકેતોની શ્રેણીઓ વચ્ચેની વાક્યરચનાત્મક અભિન્નતા
દર્શાવે છે; એટલે કે, $!A \ident !B$ ત્યારે જ જ્યારે $!A$ અને $!B$
સમાન લંબાઈ ધરાવતી સંકેતશ્રેણીઓ હોય અને દરેક સ્થાને બંનેમાં એક જ
સંકેત હોય.
\end{defn}

$\ident$ સંકેતની બંને બાજુએ જોડાણથી મળેલી શ્રેણીઓ પણ આવી શકે છે.
દાખલા તરીકે, $!A \ident (!B \lor !C)$નો અર્થ આ છે: $!A$ની
સંકેતશ્રેણી એ આરંભકૌંસ, શ્રેણી $!B$, સંકેત $\lor$, શ્રેણી $!C$ અને
અંતકૌંસને આ જ ક્રમમાં જોડવાથી મળતી શ્રેણી છે. જો એવું હોય, તો આપણે
જાણીએ છીએ કે $!A$નો પહેલો સંકેત આરંભકૌંસ છે, $!A$માં $!B$ ઉપશ્રેણી
તરીકે આવે છે (બીજા સંકેતથી શરૂ થઈને), એ ઉપશ્રેણી પછી $\lor$ આવે છે,
વગેરે.
% END OLP-0058
\endgroup


\begingroup
\makeatletter\def\ol@sectioncs{section}\makeatother

% BEGIN OLP-0059
\olfileid{pl}{syn}{pre}

\olsection{પ્રારંભિક બાબતો}
\phantomsection\label{guunit:OLP-0059}


\begin{thm}[\emph{સૂત્રો પર અનુમાનનો સિદ્ધાંત}]
  \ollabel{thm:induction}
  કોઈ ગુણધર્મ~$P$ બધાં આણ્વિક સૂત્રો માટે હોય અને નીચેની
  શરતો સંતોષતો હોય:
  \begin{enumerate}
    \item $!A$ માટે ગુણધર્મ હોય ત્યારે $\lnot !A$ માટે
      પણ હોય;
    \item $!A$ અને~$!B$ માટે ગુણધર્મ હોય ત્યારે
      $(!A \land !B)$ માટે પણ હોય;
    \item $!A$ અને~$!B$ માટે ગુણધર્મ હોય ત્યારે
      $(!A \lor !B)$ માટે પણ હોય;
    \item $!A$ અને~$!B$ માટે ગુણધર્મ હોય ત્યારે
      $(!A \lif !B)$ માટે પણ હોય;
    \item $!A$ અને~$!B$ માટે ગુણધર્મ હોય ત્યારે
      $(!A \liff !B)$ માટે પણ હોય;
  \end{enumerate}
  તો બધાં સૂત્રો માટે $P$ હોય છે.
\end{thm}

\begin{proof}
  ગુણધર્મ~$P$ ધરાવતાં બધાં સૂત્રોનો સંગ્રહ $S$ લો. સ્પષ્ટપણે
  $S \subseteq \Frm[L_0]$. $S$, \olref[fml]{defn:formulas}ની બધી
  શરતો સંતોષે છે: તેમાં બધાં આણ્વિક સૂત્રો છે અને તે
  કારકો હેઠળ સંવૃત છે. $\Frm[L_0]$ આવો નાનામાં નાનો વર્ગ છે,
  તેથી $\Frm[L_0] \subseteq S$. આમ $\Frm[L_0]=S$, અને દરેક સૂત્રને
  ગુણધર્મ~$P$ છે.
\end{proof}

\begin{prop}\ollabel{prop:balanced}
   $\Frm[L_0]$નું દરેક સૂત્ર \emph{સંતુલિત} છે, એટલે તેમાં
   ડાબાં અને જમણાં કૌંસની સંખ્યા સમાન છે.
\end{prop}

\begin{prob}
\olref[pl][syn][pre]{prop:balanced} સાબિત કરો.
\end{prob}

\begin{prop} \ollabel{prop:noinit}
સૂત્રનો કોઈ ઉચિત આરંભખંડ સૂત્ર નથી.
\end{prop}

\begin{prob}
  \olref[pl][syn][pre]{prop:noinit} સાબિત કરો.
\end{prob}

\begin{prop}[એકમાત્ર વાચનીયતા]
$\Frm[L_0]$નું દરેક સૂત્ર~$!A$ નીચેના પૈકી કોઈ એક સ્વરૂપમાં
બરાબર એક જ રીતે વિભાજિત થઈને વાંચી શકાય છે:
\begin{enumerate}
\item $\lfalse$.

\item $\ltrue$.

\item કોઈ $\Obj p_n \in \PVar$ માટે $\Obj p_n$.

\item કોઈ સૂત્ર~$!B$ માટે $\lnot !B$.

\item કોઈ સૂત્રો $!B$ અને~$!C$ માટે
$(!B \land !C)$.

\item કોઈ સૂત્રો $!B$ અને~$!C$ માટે
$(!B \lor !C)$.

\item કોઈ સૂત્રો $!B$ અને~$!C$ માટે
$(!B \lif !C)$.

\item કોઈ સૂત્રો $!B$ અને~$!C$ માટે
$(!B \liff !C)$.
\end{enumerate}
વળી, આ વિભાજિત વાંચન \emph{એકમાત્ર} છે.
\end{prop}

\begin{proof}
$!A$ પર અનુમાનથી. દાખલા તરીકે, ધારો કે $!A$નાં $(!B \lif !C)$ અને
$(!B' \lif !C')$ તરીકે બે જુદાં વાંચન છે. ત્યારે $!B$ અને $!B'$ સરખાં
હોવાં જોઈએ (નહિતર એક બીજાનો ઉચિત આરંભખંડ હોત); તેથી $!A$નાં બંને
વાંચન જુદાં હોય તો એનું કારણ માત્ર એ હોઈ શકે કે $!C$ અને $!C'$ એક જ
સંકેતશ્રેણીનાં જુદાં વાંચન છે, જે અનુમાન-પૂર્વધારણાથી અશક્ય છે.
\end{proof}


\begin{defn}[એકરૂપ પ્રતિસ્થાપન]
જો $!A$ અને $!B$ સૂત્રો હોય અને $\Obj p_i$ પ્રતિજ્ઞાપન ચલ હોય, તો $\Subst{!A}{!B}{\Obj p_i}$ એ $!A$માં $\Obj p_i$ના
દરેક આવર્તનને $!B$ના આવર્તનથી બદલવાથી મળતું પરિણામ દર્શાવે છે;
એવી જ રીતે, $\Obj p_1$, \dots,~$\Obj p_n$ને સૂત્રો $!B_1$,
\dots,~$!B_n$થી એકસાથે કરેલું પ્રતિસ્થાપન
$\SSubst{!A}{\subst{!B_1}{\Obj p_1},\dots,\subst{!B_n}{\Obj p_n}}$થી
દર્શાવાય છે.
\end{defn}

\begin{prob} નીચેનાં પાંચ સૂત્રોમાંના દરેક માટે નક્કી કરો કે
 તેને પ્રતિસ્થાપન \( \Subst{!A}{!B}{\Obj p_i} \) તરીકે લખી શકાય છે કે નહિ,
 જ્યાં \( !A \) અનુક્રમે (i) \( \Obj p_0 \); (ii) \( ( \lnot \Obj p_0 \land \Obj
 p_1) \); અને (iii) \( ( ( \lnot \Obj p_0 \lif \Obj p_1 ) \land \Obj
 p_2 ) \) છે. દરેક કિસ્સામાં સંબંધિત પ્રતિસ્થાપન જણાવો.
  \begin{enumerate}
    \item \( \Obj p_1 \)
    \item \( ( \lnot \Obj p_0 \land \Obj p_0 ) \)
    \item \( ( ( \Obj p_0 \lor \Obj p_1 ) \land \Obj p_2 )  \)
    \item \( \lnot ( ( \Obj p_0 \lif \Obj p_1 ) \land \Obj p_2 ) \)
    \item \( (( \lnot ( \Obj p_0 \lif \Obj p_1 ) \lif ( \Obj p_0 \lor \Obj p_1 )) \land \lnot ( \Obj p_0 \land \Obj p_1 )) \)
  \end{enumerate}
\end{prob}

\begin{prob}
અનુમાન વડે $\Subst{!A}{!B}{p}$ની ગાણિતિક રીતે ચુસ્ત વ્યાખ્યા આપો.
\end{prob}
% END OLP-0059
\endgroup


\begingroup
\makeatletter\def\ol@sectioncs{section}\makeatother

% BEGIN OLP-0060
\olfileid{pl}{syn}{fseq}

\olsection{રચના-શ્રેણીઓ}
\phantomsection\label{guunit:OLP-0060}


અનુમાનાત્મક વ્યાખ્યા વડે સૂત્રો વ્યાખ્યાયિત કરવા અને તેની
પૂરક રીત તરીકે અનુમાન વડે સૂત્રોના ગુણધર્મો સાબિત કરવા એ
સુઘડ અને કાર્યક્ષમ અભિગમ છે. જોકે સૂત્રોની રચનાને નીચેથી ઉપર,
ક્રમિક પગલાંમાં જોવી પણ ઉપયોગી થઈ શકે છે. અહીં આપણે \emph{રચના-શ્રેણી}
એ ખ્યાલ વડે એમ કરીશું.

\begin{defn}[સૂત્રો માટેની રચના-શ્રેણીઓ]
\ollabel{defn:fseq-frm}
ભાષા~$\Lang L_0$નાં સંકેતોની શ્રેણીઓની સાન્ત શ્રેણી
$\tuple{!A_0,\dotsc,!A_n}$ને $!A$ માટેની \emph{રચના-શ્રેણી} ત્યારે
કહીએ જ્યારે $!A \ident !A_n$ હોય અને દરેક $i\leq n$ માટે કાં તો
$!A_i$ આણ્વિક સૂત્ર હોય, અથવા એવા $j,k<i$ હોય કે નીચેનામાંથી કોઈ
એક વાત સાચી હોય:
\begin{enumerate}
    \item $!A_i \ident \lnot !A_j$.%
    \item $!A_i \ident (!A_j \land !A_k)$.%
    \item $!A_i \ident (!A_j \lor !A_k)$.%
    \item $!A_i \ident (!A_j \lif !A_k)$.%
    \item $!A_i \ident (!A_j \liff !A_k)$.%
\end{enumerate}
\end{defn}

\begin{ex}
\[
\tuple{
    \Obj p_0,
    \Obj p_1,
    (\Obj p_1 \land \Obj p_0),
    \lnot (\Obj p_1 \land \Obj p_0)
}
\]
એ $\lnot (\Obj p_1 \land \Obj p_0)$ની રચના-શ્રેણી છે; એવી જ રીતે
\[
\tuple{
    \Obj p_0,
    \Obj p_1,
    \Obj p_0,
    (\Obj p_1 \land \Obj p_0),
    (\Obj p_0 \lif \Obj p_1),
    \lnot (\Obj p_1 \land \Obj p_0)
}.
\]
પણ છે.
બીજા ઉદાહરણ પરથી દેખાય છે તેમ, રચના-શ્રેણીમાં `કચરો' હોઈ શકે છે:
એવાં સૂત્રો જે બિનજરૂરી હોય અથવા રચનામાં કોઈ ફાળો આપતાં ન હોય.
\end{ex}

\begin{prop}
\ollabel{prop:formed}
$\Frm[L_0]$ના દરેક સૂત્ર~$!A$ માટે કોઈ રચના-શ્રેણી હોય છે.
\end{prop}

\begin{proof}
ધારો કે $!A$ આણ્વિક છે. ત્યારે શ્રેણી~$\tuple{!A}$ એ $!A$ માટેની
રચના-શ્રેણી છે.
હવે ધારો કે $!B$ અને~$!C$ માટે અનુક્રમે રચના-શ્રેણીઓ
$\tuple{!B_0,\dotsc,!B_n}$ અને $\tuple{!C_0,\dotsc,!C_m}$ છે.
\begin{enumerate}
  \item જો $!A \ident \lnot !B$ હોય,
      તો $\tuple{!B_0,\dotsc,!B_n,\lnot !B_n}$
      એ $!A$ માટેની રચના-શ્રેણી છે.
  \item જો $!A \ident (!B \land !C)$ હોય,
      તો $\tuple{!B_0,\dotsc,!B_n,!C_0,\dotsc,!C_m,(!B_n \land !C_m)}$
      એ $!A$ માટેની રચના-શ્રેણી છે.
  \item જો $!A \ident (!B \lor !C)$ હોય,
      તો $\tuple{!B_0,\dotsc,!B_n,!C_0,\dotsc,!C_m,(!B_n \lor !C_m)}$
      એ $!A$ માટેની રચના-શ્રેણી છે.
  \item જો $!A \ident (!B \lif !C)$ હોય,
      તો $\tuple{!B_0,\dotsc,!B_n,!C_0,\dotsc,!C_m,(!B_n \lif !C_m)}$
      એ $!A$ માટેની રચના-શ્રેણી છે.
  \item જો $!A \ident (!B \liff !C)$ હોય,
      તો $\tuple{!B_0,\dotsc,!B_n,!C_0,\dotsc,!C_m,(!B_n \liff !C_m)}$
      એ $!A$ માટેની રચના-શ્રેણી છે.
\end{enumerate}
સૂત્રો પર અનુમાનના સિદ્ધાંત મુજબ દરેક સૂત્ર માટે
રચના-શ્રેણી હોય છે.
\end{proof}

આપણે તેનું પ્રતિલોમ પણ સાબિત કરી શકીએ છીએ. આ મહત્વનું છે, કારણ કે
તેનાથી સૂત્રો વ્યાખ્યાયિત કરવાની આપણી બંને રીતો સમતુલ્ય છે---એ બંને
સરખાં પરિણામો આપે છે---એ દેખાય છે. તેનો અર્થ એ પણ થાય છે કે
રચના-શ્રેણીની લંબાઈ પર સામાન્ય અનુમાન વાપરીને આપણે સૂત્રો વિશેનાં
પ્રમેયો સાબિત કરી શકીએ છીએ.

\begin{lem}
\ollabel{lem:fseq-init}
ધારો કે $\tuple{!A_0,\dotsc,!A_n}$ એ~$!A_n$ માટેની રચના-શ્રેણી છે
અને $k\leq n$. ત્યારે $\tuple{!A_0,\dotsc,!A_k}$ એ~$!A_k$ માટેની
રચના-શ્રેણી છે.
\end{lem}

\begin{proof}
અભ્યાસપ્રશ્ન.
\end{proof}

\begin{thm}
\ollabel{thm:fseq-frm-equiv}
$\Frm[L_0]$ એ ભાષા~$\Lang L_0$ની એવી બધી સંકેતશ્રેણીઓનો ગણ છે
જેમની રચના-શ્રેણી હોય છે.
\end{thm}

\begin{proof}
ભાષા~$\Lang L_0$ની એવી બધી સંકેતશ્રેણીઓનો ગણ $F$ લો જેમની
રચના-શ્રેણી હોય છે. \olref[pl][syn][fseq]{prop:formed}માં આપણે
$\Frm[L_0]\subseteq F$ જોયું છે; તેથી હવે આપણે તેનું પ્રતિલોમ
સાબિત કરીએ છીએ.

ધારો કે $!A$ની રચના-શ્રેણી $\tuple{!A_0,\dotsc,!A_n}$ છે. $n$ પર
પ્રબળ અનુમાનથી આપણે $!A\in\Frm[L_0]$ સાબિત કરીએ છીએ. આપણી
અનુમાન-પૂર્વધારણા એ છે કે $m<n$ લંબાઈની રચના-શ્રેણી ધરાવતી દરેક
સંકેતશ્રેણી $\Frm[L_0]$માં છે. રચના-શ્રેણીની વ્યાખ્યા મુજબ કાં તો
$!A_n$ આણ્વિક છે, અથવા એવા $j,k<n$ હોવા જોઈએ કે નીચેનામાંથી કોઈ એક
કિસ્સો બને:
\begin{enumerate}
    \item $!A_n \ident \lnot !A_j$.%
    \item $!A_n \ident (!A_j \land !A_k)$.%
    \item $!A_n \ident (!A_j \lor !A_k)$.%
    \item $!A_n \ident (!A_j \lif !A_k)$.%
    \item $!A_n \ident (!A_j \liff !A_k)$.%
\end{enumerate}
હવે આપણે કિસ્સાવાર વિચારીએ. જો $!A_n$ આણ્વિક હોય, તો
$!A_n\in\Frm[L_0]$. હવે તેના બદલે ધારો કે $!A\ident
(!A_j\land !A_k)$. \olref[pl][syn][fseq]{lem:fseq-init} મુજબ
$\tuple{!A_0,\dotsc,!A_j}$ અને $\tuple{!A_0,\dotsc,!A_k}$ અનુક્રમે
$!A_j$ અને~$!A_k$ માટેની રચના-શ્રેણીઓ છે. આ બંને $!A$ની
રચના-શ્રેણીની ઉચિત આરંભ-ઉપશ્રેણીઓ હોવાથી તેમની લંબાઈ~$n$થી ઓછી છે.
તેથી અનુમાન-પૂર્વધારણા મુજબ $!A_j$ અને~$!A_k$, $\Frm[L_0]$માં છે;
અને આમ સૂત્રની વ્યાખ્યા મુજબ $(!A_j\land !A_k)$ પણ તેમાં છે.
બીજા કિસ્સા સમાંતર દલીલથી મળે છે.
\end{proof}
\sourcecorrection{OLPL-002}{મૂળની \texttt{formation-sequences.tex}માં
આ કિસ્સાની પૂર્વધારણા $!A\equiv(!A_j\land !A_k)$ લખાઈ છે. અહીં
તાર્કિક સમતુલ્યતાની નહિ પણ સંકેતશ્રેણીની અભિન્નતાની જરૂર છે;
વ્યાખ્યામાં અને આ જ સાબિતીના અગાઉના ખંડમાં વપરાયેલો
$\ident$ સંકેત અહીં પુનઃસ્થાપિત કર્યો છે.}
% END OLP-0060
\endgroup


\begingroup
\makeatletter\def\ol@sectioncs{section}\makeatother

% BEGIN OLP-0061
\olfileid{pl}{syn}{val}

\olsection{સત્યમૂલ્ય-નિયુક્તિઓ અને સંતોષ}
\phantomsection\label{guunit:OLP-0061}


\begin{defn}[સત્યમૂલ્ય-નિયુક્તિઓ]
$\{\True,\False\}$ને ``સત્ય'' અને ``અસત્ય'' એ બે સત્યમૂલ્યોનો ગણ
માનો. $\Lang{L_0}$ માટેની \emph{સત્યમૂલ્ય-નિયુક્તિ} એ ભાષાનાં
પ્રતિજ્ઞાપન ચલોને $\True$ અથવા $\False$ નિયુક્ત કરતું
વિધેય~$\pAssign{v}$ છે; એટલે કે, $\pAssign{v}\colon
\PVar\to\{\True,\False\}$.
\end{defn}

\begin{defn}
  સત્યમૂલ્ય-નિયુક્તિ~$\pAssign{v}$ આપેલી હોય, ત્યારે મૂલ્યાંકન વિધેય
  $\pValue{v}\colon\Frm[L_0]\to\{\True,\False\}$ને નીચે પ્રમાણે
  અનુમાનાત્મક રીતે વ્યાખ્યાયિત કરો:
  \begin{align*}
    \pValue{v}(\lfalse) & = \False; \\
    \pValue{v}(\ltrue) & = \True; \\
    \pValue{v}(\Obj p_n) & = \pAssign{v}(\Obj p_n); \\
    \pValue{v}(\lnot !A) & = \begin{cases}
        \True & \text{જો } \pValue{v}(!A) = \False;\\
        \False & \text{અન્યથા.}
      \end{cases} \\ 
    \pValue{v}(!A \land !B) & = \begin{cases}
        \True &
        \text{જો $\pValue{v}(!A) = \True$ અને $\pValue{v}(!B) = \True$;}\\
        \False &
        \text{જો $\pValue{v}(!A) = \False$ અથવા $\pValue{v}(!B) = \False$}.
      \end{cases}\\
    \pValue{v}(!A \lor !B) & = \begin{cases}
        \True &
        \text{જો $\pValue{v}(!A) = \True$ અથવા $\pValue{v}(!B) = \True$;}\\
        \False &
        \text{જો $\pValue{v}(!A) = \False$ અને $\pValue{v}(!B) = \False$}.
      \end{cases}\\
    \pValue{v}(!A \lif !B) & = \begin{cases}
        \True &
        \text{જો $\pValue{v}(!A) = \False$ અથવા $\pValue{v}(!B) = \True$;}\\
        \False &
        \text{જો $\pValue{v}(!A) = \True$ અને $\pValue{v}(!B) = \False$}.
      \end{cases}\\
    \pValue{v}(!A \liff !B) & = \begin{cases}
        \True &
        \text{જો $\pValue{v}(!A) = \pValue{v}(!B)$;}\\
        \False &
        \text{જો $\pValue{v}(!A) \neq \pValue{v}(!B)$}.
      \end{cases}
\end{align*}
\end{defn}

\begin{explain}
આ ખંડો નીચેની સત્યાર્થતા સારણીઓને અનુરૂપ છે:
\begin{center}
  \begin{tabular}{|c||c|} \hline
    $!A$ & $ \lnot !A$ \\
    \hline \hline
    $\True$ & $\False$ \\
    $\False$ & $\True$ \\
    \hline
  \end{tabular}
  \begin{tabular}{|cc||c|} \hline
    $!A$ & $!B$ & $!A \land !B$ \\
    \hline \hline
    $\True$ & $\True$ & $\True$ \\
    $\True$ & $\False$ & $\False$ \\
    $\False$ & $\True$ & $\False$ \\
    $\False$ & $\False$ & $\False$ \\
    \hline
  \end{tabular}
  \begin{tabular}{|cc||c|} \hline
    $!A$ & $!B$ & $!A \lor !B$ \\
    \hline \hline
    $\True$ & $\True$ & $\True$ \\
    $\True$ & $\False$ & $\True$ \\
    $\False$ & $\True$ & $\True$ \\
    $\False$ & $\False$ & $\False$ \\
    \hline
  \end{tabular}
%
\\[1em]
  \begin{tabular}{|cc||c|} \hline
    $!A$ & $!B$ & $!A \lif !B$ \\
    \hline \hline
    $\True$ & $\True$ & $\True$ \\
    $\True$ & $\False$ & $\False$ \\
    $\False$ & $\True$ & $\True$ \\
    $\False$ & $\False$ & $\True$ \\
    \hline
  \end{tabular}
  \begin{tabular}{|cc||c|} \hline
    $!A$ & $!B$ & $!A \liff !B$ \\
    \hline \hline
    $\True$ & $\True$ & $\True$ \\
    $\True$ & $\False$ & $\False$ \\
    $\False$ & $\True$ & $\False$ \\
    $\False$ & $\False$ & $\True$ \\
    \hline
  \end{tabular}
\end{center}
\end{explain}

\begin{prob}
$\Lang{L_0}$માં ત્રિસ્થાની સંયોજક $\diamondsuit$ ઉમેરવાનું વિચારો,
જેનું મૂલ્યાંકન આ પ્રમાણે આપેલું છે:
\begin{gather*}
  \pValue{v}(\diamondsuit( !A , !B , !C ) ) = \begin{cases}
    \pValue{v}( !B ) &
    \text{જો $\pValue{v}(!A) = \True$;}\\
    \pValue{v}( !C ) &
    \text{જો $\pValue{v}(!A) = \False $}.
  \end{cases}
\end{gather*}
આ સંયોજકની સત્યાર્થતા સારણી લખો.
\end{prob}

\begin{thm}[સ્થાનિક નિર્ધારણ]
\ollabel{thm:LocalDetermination} ધારો કે $\pAssign{v_1}$ અને
$\pAssign{v_2}$ એ સત્યમૂલ્ય-નિયુક્તિઓ છે અને $!A$માં આવતાં
પ્રતિજ્ઞાપન ચલો પર સંમત થાય છે; એટલે કે, કોઈ
સૂત્ર~$!A$માં $\Obj p_n$ આવે ત્યારે
$\pAssign{v_1}(\Obj p_n)=\pAssign{v_2}(\Obj p_n)$. ત્યારે
$\pValue{v_1}$ અને $\pValue{v_2}$ પણ~$!A$ પર સંમત થાય છે; એટલે કે,
$\pValue{v_1}(!A)=\pValue{v_2}(!A)$.
\end{thm}

\begin{proof}
$!A$ પર અનુમાનથી.
\end{proof}

\begin{defn}[સંતોષ]
\ollabel{defn:satisfaction} સત્યમૂલ્ય-નિયુક્તિ~$\pAssign{v}$ વડે
  સૂત્ર~$!A$ના \emph{સંતોષ}, $\pSat{v}{!A}$,નો ખ્યાલ આપણે
  નીચે પ્રમાણે અનુમાનાત્મક રીતે વ્યાખ્યાયિત કરી શકીએ છીએ.
  (``$\pSat{v}{!A}$ નથી'' એ કહેવા આપણે $\pSat/{v}{!A}$ લખીએ છીએ.)
\begin{enumerate}
\item %
  \indcase{!A}{\lfalse}{$\pSat/{v}{\indfrm}$.}

\item %
  \indcase{!A}{\ltrue}{$\pSat{v}{\indfrm}$.}

\item \indcase{!A}{\Obj p_i}{$\pSat{v}{\indfrm}$ ત્યારે જ જ્યારે
  $\pAssign{v}(\Obj p_i) = \True$.}

\item %
  \indcase{!A}{\lnot !B}{$\pSat{v}{\indfrm}$ ત્યારે જ જ્યારે
    $\pSat/{v}{!B}$.}

\item %
  \indcase{!A}{(!B \land !C)}{$\pSat{v}{\indfrm}$ ત્યારે જ જ્યારે
    $\pSat{v}{!B}$ અને $\pSat{v}{!C}$.}

\item %
  \indcase{!A}{(!B \lor !C)}{$\pSat{v}{\indfrm}$ ત્યારે જ જ્યારે
    $\pSat{v}{!B}$ અથવા $\pSat{v}{!C}$ (અથવા બંને).}

\item %
  \indcase{!A}{(!B \lif !C)}{$\pSat{v}{\indfrm}$ ત્યારે જ જ્યારે
    $\pSat/{v}{!B}$ અથવા $\pSat{v}{!C}$ (અથવા બંને).}

\item %
  \indcase{!A}{(!B \liff !C)}{$\pSat{v}{\indfrm}$ ત્યારે જ જ્યારે
    $\pSat{v}{!B}$ અને $\pSat{v}{!C}$ બંને, અથવા $\pSat{v}{!B}$ અને
    $\pSat{v}{!C}$ એ બેમાંથી એકેય નહિ.}
\end{enumerate}
જો $\Gamma$ એ સૂત્રોનો ગણ હોય, તો દરેક~$!A\in\Gamma$ માટે
$\pSat{v}{!A}$ હોય ત્યારે જ $\pSat{v}{\Gamma}$.
\end{defn}

\begin{prop}\ollabel{prop:sat-value}
  $\pSat{v}{!A}$ ત્યારે જ જ્યારે $\pValue{v}(!A)=\True$.
\end{prop}

\begin{proof}
  $!A$ પર અનુમાનથી.
\end{proof}

\begin{prob}
  \olref[pl][syn][val]{prop:sat-value} સાબિત કરો.
\end{prob}
% END OLP-0061
\endgroup


\begingroup
\makeatletter\def\ol@sectioncs{section}\makeatother

% BEGIN OLP-0062
\olfileid{pl}{syn}{sem}

\olsection{અર્થવિચારી ખ્યાલો}
\phantomsection\label{guunit:OLP-0062}


આપણે નીચેના અર્થવિચારી ખ્યાલો વ્યાખ્યાયિત કરીએ છીએ:

\begin{defn}
\begin{enumerate}
\item કોઈ~$\pAssign{v}$ માટે $\pSat{v}{!A}$ હોય, તો
  સૂત્ર~$!A$ \emph{સંતોષ્ય} છે; કોઈ પણ $\pAssign{v}$ માટે
  $\pSat{v}{!A}$ ન હોય, તો તે \emph{અસંતોષ્ય} છે;
\item બધાં સત્યમૂલ્ય-નિયુક્તિઓ~$\pAssign{v}$ માટે $\pSat{v}{!A}$ હોય, તો
  સૂત્ર~$!A$ \emph{પુનરુક્તિ} છે;
\item સૂત્ર~$!A$ સંતોષ્ય હોય પણ પુનરુક્તિ ન હોય, તો તે
  \emph{નિવાર્ય} છે;
\item જો $\Gamma$ એ સૂત્રોનો ગણ હોય, તો દરેક એવા
  સત્યમૂલ્ય-નિયુક્તિ~$\pAssign{v}$ માટે કે જેના માટે $\pSat{v}{\Gamma}$,
  $\pSat{v}{!A}$ પણ હોય ત્યારે અને માત્ર ત્યારે $\Gamma\Entails !A$
  (``$\Gamma$, $!A$ને ફલિત કરે છે'').
\item જો $\Gamma$ એ સૂત્રોનો ગણ હોય, તો એવા કોઈ
  સત્યમૂલ્ય-નિયુક્તિ~$\pAssign{v}$ માટે $\pSat{v}{\Gamma}$ હોય ત્યારે
  $\Gamma$ \emph{સંતોષ્ય} છે; અન્યથા $\Gamma$ \emph{અસંતોષ્ય} છે.
\end{enumerate}
\end{defn}

\begin{prob}
 નીચેનાં ચાર સૂત્રોમાંનું દરેક (a)~સંતોષ્ય, (b)~પુનરુક્તિ અને
 (c)~નિવાર્ય છે કે નહિ તે નક્કી કરો.
\begin{enumerate}
  \item \( ( \Obj p_0 \lif ( \lnot \Obj p_1 \lif \lnot \Obj p_0 ) ) \).
  \item \( ( ( \Obj p_0 \land \lnot \Obj p_1 ) \lif ( \lnot \Obj p_0 \land \Obj p_2 )) \liff ( ( \Obj p_2 \lif \Obj p_0 ) \lif ( \Obj p_0 \lif \Obj p_1 )) \).
  \item \( ( \Obj p_0 \liff \Obj p_1 ) \lif ( \Obj p_2 \liff \lnot \Obj p_1 ) \).
  \item \( (( \Obj p_0 \liff ( \lnot \Obj p_1 \land \Obj p_2 )) \lor ( \Obj p_2 \lif ( \Obj p_0 \liff \Obj p_1 ))) \).
\end{enumerate}
\end{prob}

\begin{prop}
\ollabel{prop:semanticalfacts}
\begin{enumerate}
\item $!A$ પુનરુક્તિ હોય ત્યારે અને માત્ર ત્યારે
  $\emptyset \Entails !A$;
\item જો $\Gamma \Entails !A$ અને $\Gamma \Entails !A \lif !B$, તો
  $\Gamma \Entails !B$;
\item જો $\Gamma$ સંતોષ્ય હોય, તો $\Gamma$નો દરેક સાન્ત ઉપગણ પણ
  સંતોષ્ય છે;
\item \ollabel{def:monotonicity} એકદિશવર્ધિતા: જો
  $\Gamma \subseteq \Delta$ અને $\Gamma \Entails !A$, તો
  $\Delta \Entails !A$ પણ;
\item \ollabel{def:Cut} પરંપરિતતા: જો $\Gamma \Entails !A$ અને
  $\Delta \cup \{ !A\} \Entails !B$, તો $\Gamma \cup \Delta \Entails
  !B$.
\end{enumerate}
\end{prop}

\begin{proof}
અભ્યાસપ્રશ્ન.
\end{proof}

\begin{prob}
\olref[pl][syn][sem]{prop:semanticalfacts} સાબિત કરો.
\end{prob}

\begin{prop}\ollabel{prop:entails-unsat}
  $\Gamma \Entails !A$ ત્યારે અને માત્ર ત્યારે જ્યારે
  $\Gamma \cup \{\lnot !A\}$ અસંતોષ્ય હોય.
\end{prop}

\begin{proof}
અભ્યાસપ્રશ્ન.
\end{proof}

\begin{prob}
\olref[pl][syn][sem]{prop:entails-unsat} સાબિત કરો.
\end{prob}

\begin{thm}[અર્થાનુસારી નિગમન પ્રમેય]
  \ollabel{thm:sem-deduction} $\Gamma \Entails !A \lif !B$ ત્યારે અને માત્ર
  ત્યારે જ્યારે $\Gamma \cup \{!A\} \Entails !B$.
\end{thm}

\begin{proof}
અભ્યાસપ્રશ્ન.
\end{proof}

\begin{prob}
\olref[pl][syn][sem]{thm:sem-deduction} સાબિત કરો.
\end{prob}
% END OLP-0062
\endgroup


\OLEndChapterHook


\begin{editorial}
આગળનો સંબંધિત અંશ મૂળના સક્રિય આયાતક્રમમાં નથી. તેને અહીં સ્પષ્ટ વૈકલ્પિક સંદર્ભમાં જાળવ્યો છે.
\end{editorial}
\begingroup
\makeatletter\def\ol@sectioncs{section}\makeatother

% BEGIN OLP-0714
\olfileid{pl}{prp}{com}

\olsection{પ્રતિજ્ઞાપન તર્કશાસ્ત્રની પૂર્ણતા}
\phantomsection\label{guunit:OLP-0714}


\begin{defn}
\ollabel{def:MaxCon}
સૂત્રનો ગણ~$\Gamma$ \emph{મહત્તમ સુસંગત} કહેવાય,
જો તે સુસંગત હોય અને કોઈ સુસંગત ગણ~$\Delta$ માટે
$\Gamma \subseteq \Delta$ હોય ત્યારે $\Gamma = \Delta$ થાય.
\end{defn}

\begin{lem}[સત્ય સહાયક પ્રમેય]
  \ollabel{lem:truth}
$\Gamma$ મહત્તમ સુસંગત હોય, તો:
\begin{enumerate}
\item \ollabel{lem:truth-1} $!A \in \Gamma$ જો અને ફક્ત જો
  $\lnot !A \notin \Gamma$;
\item \ollabel{lem:truth-2} $\Gamma \Proves !A$ જો અને ફક્ત જો
  $!A \in \Gamma$;
\item \ollabel{lem:truth-3} $!A \lif !B \in \Gamma$ જો અને ફક્ત જો
  કાં તો $!A \notin \Gamma$ અથવા $!B \in \Gamma$.
\end{enumerate}
\end{lem}

\begin{proof}
\begin{enumerate}
\item ધારો કે $!A \in \Gamma$ હોય. જો $\lnot !A \in \Gamma$
  પણ હોય, તો $\Gamma$ અસુસંગત થાય.
  બીજી શક્યતા એ છે કે $!A$ કે $\lnot!A$ બેમાંથી એકેય
  $\Gamma$માં ન હોય. ત્યારે
  \olref[fol][axd][prv]{prop:provability-exhaustive} પ્રમાણે
  $\Gamma\cup\{!A\}$ અને $\Gamma \cup \{\lnot !A\}$માંથી
  એક સુસંગત હોય, એટલે $\Gamma$ મહત્તમ ન હોય.
  આ બંને અવલોકનો ઊલટી દિશા પણ આપે છે.
\item જો $!A \in \Gamma$ હોય, તો દેખીતી રીતે $\Gamma \Proves !A$.
  ઊલટી તરફ, ધારો કે $\Gamma \Proves !A$ હોય.
  જો $!A \notin \Gamma$ હોય, તો \olref{lem:truth-1}
  પ્રમાણે $\lnot !A \in \Gamma$ થાય.
  આમ $\Gamma \Proves !A$ અને $\Gamma \Proves \lnot !A$,
  જે સુસંગતતા સાથે વિરોધાભાસ આપે છે.
\item કસરત.
\end{enumerate}
\sourcecorrection{OLMD-344}{મૂળના બે જૂના સંદર્ભ એકેય વિસ્તરણ સુસંગત ન હોય તો મૂળ ગણ અસુસંગત થાય તે પરિણામને નિર્દેશે છે, પરંતુ આપેલા નામવાળો ઉલ્લેખ સ્થિર સ્રોતમાં મળતો નથી. તેમની જગ્યાએ હાલની સ્વયંસિદ્ધ-નિગમન રચનાના અસુસંગતતાનાં બંને વિસ્તરણ અંગેના પરિણામનો સ્પષ્ટ સંદર્ભ આપ્યો છે. આગળના નિષ્પન્નતા-અસુસંગતતા સંદર્ભને પણ તેના સાચા વિભાગ સુધી વિસ્તૃત કર્યો છે.}
\end{proof}

\olref{lem:truth}\olref{lem:truth-2}ની ડાબેથી જમણી દિશા
બતાવે છે કે $\Gamma$ \emph{નિગમન હેઠળ બંધ} છે:
દરેક $!A$ માટે, જો $\Gamma \Proves !A$ હોય, તો $!A
\in \Gamma$ થાય.

\begin{prob}
  \olref{lem:truth}ની સાબિતી પૂર્ણ કરો.
\end{prob}

\begin{thm}[પૂર્ણતા]
\ollabel{thm:completeness}
જો $\Gamma$ સુસંગત હોય, તો $\Gamma$ સંતોષ્ય છે.
\end{thm}

\begin{proof}
ભાષાનાં બધાં સૂત્રની સંપૂર્ણ સૂચિ $!A_0$, $!A_1$,
\dots લો. સૂત્રના ગણોનો વધતો ક્રમ
$\Gamma_0$, $\Gamma_1$, \dots નીચે મુજબ પુનરાવર્તનથી વ્યાખ્યાયિત કરીએ:
\begin{align*}
\Gamma_0 & = \Gamma\\
\Gamma_{n+1} &  =
\begin{cases}
  \Gamma_n \cup \{!A_n\} & \text{જો $\Gamma_n \cup \{!A_n\}$ સુસંગત હોય;}\\
  \Gamma_n \cup \{\lnot !A_n\} & \text{અન્યથા}.
\end{cases}
\end{align*}
પછી આ વ્યાખ્યા આપીએ:
\[
\Gamma^* = \bigcup_{0\le n}\Gamma_n.
\]
હવે સાબિતીને આગળ વધારવા ક્રમશઃ નીચેની હકીકતો સ્થાપિત કરવાની છે:
\begin{enumerate}
\item દરેક $n$ માટે ગણ~$\Gamma_n$ સુસંગત છે:
  \olref[fol][axd][prv]{prop:provability-exhaustive} વાપરી
  $n$ પર આગમન કરો;
\item $\Gamma^*$ સુસંગત છે;
\item $\Gamma^*$ મહત્તમ છે.
\end{enumerate}
પછી સત્યમૂલ્ય-નિયુક્તિ~$v$ એવી વ્યાખ્યાયિત કરીએ કે
$v(p_i) = \True$ જો અને ફક્ત જો $p_i \in \Gamma^*$.
$!A$ પર આગમનથી બતાવવાનું છે કે વધુ જટિલ વાક્ય
માટે પણ $\Gamma^*$નું સભ્યપદ અને $v$ અનુસારનું સત્ય એક જ છે:
$\overline{v}(!A) = \True$ જો અને ફક્ત જો $!A \in \Gamma^*$.
ખાસ કરીને $v$ એ $\Gamma^*$ને સંતોષે છે;
$\Gamma \subseteq \Gamma^*$ હોવાથી $\Gamma$ પણ સંતોષ્ય છે,
જે જોઈએ છે.
\end{proof}

\begin{cor}
જો $\Gamma \Entails !A$ હોય, તો $\Gamma \Proves !A$.
\end{cor}

\begin{proof}
જો $\Gamma\Proves/ !A$ હોય, તો
\olref[fol][axd][prv]{prop:prov-incons} પ્રમાણે
$\Gamma \cup \{\lnot !A\}$ સુસંગત છે.
પૂર્ણતા પ્રમેયથી $\Gamma \cup
\{\lnot !A\}$ સંતોષ્ય છે, એટલે $\Gamma \Entails/ !A$.
\end{proof}

\begin{prop}[સઘનતા પ્રમેય]
\ollabel{prop:compactness}
$\Gamma$ સંતોષ્ય છે જો અને ફક્ત જો $\Gamma$નો દરેક
\emph{સાંત} ઉપગણ~$\Gamma_0$ સંતોષ્ય હોય.
\end{prop}

\begin{proof}
$\Gamma$ અસંતોષ્ય છે જો અને ફક્ત જો તે અસુસંગત હોય;
જો અને ફક્ત જો $\Gamma$નો કોઈ સાંત ઉપગણ~$\Gamma_0$
અસુસંગત હોય; જો અને ફક્ત જો $\Gamma$નો કોઈ સાંત
ઉપગણ~$\Gamma_0$ અસંતોષ્ય હોય.
\end{proof}

\begin{cor}
$\Gamma \Entails !A$ જો અને ફક્ત જો $\Gamma$ના
કોઈ સાંત ઉપગણ~$\Gamma_0$ માટે $\Gamma_0 \Entails !A$.
\end{cor}

\begin{editorial}
આ પઠન-રચનામાં ફરી વપરાયેલા નિષ્પન્નતા-ગુણધર્મોનો પાઠ પ્રથમ-ક્રમના સ્વયંસિદ્ધિ-આધારિત નિગમનના ભાગમાં એક જ વાર છે. આ વિભાગની કડીઓ ત્યાંના તે જ પરિણામોના વિધાનાત્મક કિસ્સાઓ સુધી જાય છે.
\end{editorial}
% END OLP-0714
\endgroup

\begin{editorial}
આગળનો સંબંધિત અંશ મૂળના સક્રિય આયાતક્રમમાં નથી. તેને અહીં સ્પષ્ટ વૈકલ્પિક સંદર્ભમાં જાળવ્યો છે.
\end{editorial}
\begingroup
\makeatletter\def\ol@sectioncs{section}\makeatother

% BEGIN OLP-0715
\olfileid{pl}{prp}{snd}

\olsection{પ્રતિજ્ઞાપન તર્કશાસ્ત્રની યથાર્થતા}
\phantomsection\label{guunit:OLP-0715}


\begin{thm}[યથાર્થતા]
\ollabel{thm:soundness}
જો $\Gamma\Proves!A$ હોય, તો $\Gamma \Entails !A$.
\end{thm}

\begin{proof}
નિષ્પન્ન થતી પ્રતિજ્ઞાપનાઓ પર આગમન કરીએ.
જો $!A$ સ્વયંસિદ્ધ હોય, તો $\Entails
!A$, એટલે $\Gamma \Entails!A$ પણ થાય.
તે જ રીતે, જો $!A \in \Gamma$ હોય, તો $\Gamma \Entails!A$.
આગમનના પગલામાં ધારો કે $!A$ એ $!C$ અને $!C\lif!A$માંથી
\emph{મોડસ પોનેન્સ} વડે મળ્યું હોય.
આગમનની ધારણા મુજબ $!C\lif !A$ અને $!C$ માટે પ્રમેય સાચું છે.
\olref[pl][syn][sem]{prop:semanticalfacts}નો બીજો મુદ્દો
ઇચ્છિત પરિણામ આપે છે.
\sourcecorrection{OLMD-345}{મૂળ અહીં અર્થવિજ્ઞાનના પરિણામનો ત્રીજો મુદ્દો કહે છે. સ્થિર મૂળમાં મોડસ પોનેન્સને અનુરૂપ સત્ય-સંરક્ષણ બીજો મુદ્દો છે; ત્રીજો મુદ્દો સાંત ઉપગણોની સંતોષ્યતા અંગે છે. ક્રમસંખ્યા અને સ્રોતના સાચા ભાગ-અધ્યાય સુધીનો સંદર્ભ સુધાર્યા છે.}
\end{proof}

\begin{cor}
જો $\Gamma$ સંતોષ્ય હોય, તો $\Gamma$ સુસંગત છે.
આથી પ્રતિજ્ઞાપન તર્કશાસ્ત્ર સુસંગત છે.
\end{cor}
% END OLP-0715
\endgroup
% END OLP-0056
\endgroup


\tagfalse{FOL}

\noindent\hyperref[guunit:OLP-0063]{સંબંધિત સામગ્રી જુઓ (OLP-0063).}\par


%
  \noindent\hyperref[guunit:OLP-0069]{સંબંધિત સામગ્રી જુઓ (OLP-0069).}\par


%
  \noindent\hyperref[guunit:OLP-0084]{સંબંધિત સામગ્રી જુઓ (OLP-0084).}\par


%
  \noindent\hyperref[guunit:OLP-0098]{સંબંધિત સામગ્રી જુઓ (OLP-0098).}\par


%
  \noindent\hyperref[guunit:OLP-0112]{સંબંધિત સામગ્રી જુઓ (OLP-0112).}\par


\noindent\hyperref[guunit:OLP-0126]{સંબંધિત સામગ્રી જુઓ (OLP-0126).}\par


\tagtrue{FOL}

\OLEndPartHook
% END OLP-0055
